% !TeX root = ../main.tex
\chapter{METODOLOGI PENELITIAN}
% Kerangka saja; uraian, diagram, tabel, dan rumus ditambahkan pada tahap berikutnya.

\section{Data Penelitian}

\subsection{Data Rekaman Suara Burung}
% Sumber Xeno-Canto, 20 label spesies, metadata, dan unit rekaman aktif.
% Bedakan persebaran spesies dari negara asal tiap rekaman.

\subsection{Data Derau Lapangan ITERA}
% Lima lokasi, proses perekaman, dan peran bank derau pada eksperimen.
% Bedakan pool pencampuran dari derau yang digunakan sebagai kontrol di luar target.

\subsection{Kriteria Pemilihan dan Pembagian Data}
% Galeri, kueri, kalibrasi, dan kontrol di luar target mengikuti manifes aktif.
% Jelaskan pemisahan perekam dan pencegahan kebocoran data.

\section{Diagram Alir Penelitian}
% Diagram alir akan disusun berdasarkan data, tahapan, dan eksperimen yang benar-benar dilakukan.
% Bagian ini masih berupa kerangka; gambar dan penjelasan belum ditambahkan.

\section{Tahapan Penelitian}

\subsection{Rancangan Eksperimen}
% Pendekatan kuantitatif, variabel, kondisi acuan, kontrol, dan kelompok eksperimen.
% Bedakan rancangan awal dari analisis eksploratif setelah hasil tersedia.

\subsection{Perangkat Penelitian}
% Perangkat perekaman serta perangkat keras/lunak sesuai bukti yang tersedia.

\subsection{Prapemrosesan Audio}
% Penyeragaman kanal, laju sampel, durasi, dan tingkat sinyal sesuai kode aktif.

\subsection{Ekstraksi Representasi Audio}
% MFCC, PANNs CNN14, BirdNET, serta vektor acak berbasis identitas sebagai kontrol.
% Bobot model pralatih tetap; MFCC merupakan fitur hitungan.

\subsection{Pencarian Rekaman Berdasarkan Kemiripan Suara}
% Penyusunan galeri tetap, kemiripan kosinus, dan pengurutan hasil pencarian.

\subsection{Evaluasi Kualitas Pencarian}
% mAP@10 sebagai metrik utama; precision@10, recall@10, MRR, dan Top-1 sebagai pendukung.
% Penilaian relevansi mengikuti label spesies dan hasil pencarian untuk setiap kueri.

\subsection{Pengujian Ketahanan pada Berbagai Tingkat Derau}
% Kondisi tanpa derau tambahan dan pencampuran aditif pada empat SNR.
% Bandingkan kueri yang sama dengan galeri tetap menggunakan kualitas dan retensi.

\subsection{Kalibrasi Ambang dan Pengujian Penerimaan Kueri}
% Kalibrasi terpisah tanpa penambahan derau; nilai ambang dipertahankan saat uji.
% Pisahkan penerimaan target, penerimaan keliru kontrol, dan ketepatan hasil pertama.

\subsection{Pengujian Profil Derau Lima Lokasi}
% Bandingkan lima profil pada kueri dan SNR yang sama dengan empat ulangan.
% Laporkan ayunan absolut dan relatif, serta sensitivitas tanpa Kebun Raya sebagai analisis eksploratif.

\subsection{Analisis Kegagalan dan Audit Pendengaran}
% Margin terhadap ambang, ciri spektral, dan audit 120 kasus pilihan oleh satu pendengar.
% Analisis deskriptif; bukan bukti penyebab kegagalan atau validasi spesies independen.
% Ciri akustik tersedia pada E2 yang direkonstruksi, tidak disimpulkan untuk E4.

\subsection{Analisis Statistik dan Interpretasi}
% Analisis bootstrap berpasangan pada kueri untuk sensitivitas kestabilan profil.
% Jelaskan unit analisis, interval kepercayaan, keterlacakan, dan batas interpretasi.

\section{Jadwal Penelitian}
% Jadwal mengikuti kegiatan dan catatan pelaksanaan yang dapat ditelusuri.
